\begin{center}
\textsc{Gauss} an \textsc{Gerling}.\quad Göttingen, 7.~November 1822.
\end{center}

\par\noindent\hspace*{\parindent}\mbox{.\ .\ .\ .\ .\ .}\quad Bei sehr flacher Incidenz, wo auch zuletzt die Möglichkeit der
Lenkung aufhört, habe ich immer mit dem \textls[120]{herrlichsten} Erfolg doppelte
Reflexion anwenden lassen, indem nicht die Sonne selbst, sondern ein auf der
Erde an schicklicher Stelle aufgestellter Handspiegel den Heliotrop speiste.
Ich empfehle Ihnen diesen Kunstgriff zur Nachahmung.\quad\mbox{.\ .\ .\ .\ .\ .}

\par\noindent\hspace*{\parindent}In einem bergigen Lande ist es eine Lust zu messen, desto grösser, je
grösser die Entfernungen sind, die beim Gebrauch der Heliotropen gar keine
Grenzen haben. Mein \textls[120]{grosser} Spiegel, 1 Quadratfuss, war sogar durch den
dichten Moorbrandsqualm, der mehrere Quadratmeilen bedeckte, von
Lichtenberg nach Falkenberg, fast 12 Meilen, durchgedrungen, welcher Qualm freilich
für die winzigen Heliotropspiegel zu dicht gewesen war.\quad\mbox{.\ .\ .\ .\ .\ .}
